%----------------------------------------------------------------------------------------------------
% START OF CHAPTER
%----------------------------------------------------------------------------------------------------
\part{Core Verbs} \label{M2-Core Verbs}

\chapter{\iti{Essere} and \iti{Avere}} \label{M2-L1}

\section{\iti{Essere} — to be}
\begin{table}[ht]
    \centering
    \begin{tabular}{r l}
      Person & Form \\
      \hline
      io      & \iti{sono}  \\
      tu      & \iti{sei}   \\
      lui/lei & \iti{è}     \\
      noi     & \iti{siamo} \\
      voi     & \iti{siete} \\
      loro    & \iti{sono} \\
    \end{tabular}
\end{table}

\section{\iti{Avere} — to have}
\begin{table}[ht]
    \centering
    \begin{tabular}{r l}
        Person & Form \\
        \hline
        io      & \iti{ho}  \\
        tu      & \iti{hai}   \\
        lui/lei & \iti{ha}     \\
        noi     & \iti{abbiamo} \\
        voi     & \iti{avete} \\
        loro    & \iti{hanno} \\
    \end{tabular}
\end{table}

The H in \iti{ho, hai, ha, hanno} is \textbf{always} silent.

\section{\iti{Avere} for feelings}

\textbf{!! not} \iti{essere}!

\begin{table}[ht]
    \centering
    \begin{tabular}{r l}
        Italian & English \\
        \hline
        \iti{ho fame}   & I am hungry  \\
        \iti{ho sete}   & I am thirsty \\
        \iti{ho freddo} & I am cold    \\
        \iti{ho caldo}  & I am hot     \\
        \iti{ho paura}  & I am afraid  \\
        \iti{ho sonno}  & I am sleepy  \\
        \iti{ho ... anni} & I am ... years old \\
    \end{tabular}
\end{table}

\textbf{Alert!!} \iti{Sono caldo} means something very different from \iti{ho caldo} — always use \iti{avere} for these!


\chapter{Regular Verbs — Present Tense} \label{M2-L2}

\section{The Four Anchor Endings (all verb families)}
\begin{itemize}[nosep]
    \item \iti{io} → \itb{-o}
    \item \iti{tu} → \itb{-i}
    \item \iti{noi} → \itb{-iamo}
    \item \iti{loro} → \itb{-no}
\end{itemize}

\section{-ARE Verbs (\iti{parlare})}
\iti{parlo, parli, parla, parliamo, parlate, parlano}

Common: \iti{mangiare, lavorare, abitare, ascoltare, comprare, camminare}

\section{-ERE Verbs (\iti{leggere})}
\iti{leggo, leggi, legge, leggiamo, leggete, leggono}

Common: \iti{vedere, scrivere, prendere, vivere, credere, chiudere}

\section{-IRE Verbs (\iti{dormire})}
\iti{dormo, dormi, dorme, dormiamo, dormite, dormono}

Common: \iti{partire, sentire, aprire, seguire}

\section{-IRE Verbs with -ISC- (\iti{capire})}
\iti{capisco, capisci, capisce, capiamo, capite, capiscono}

Note! \iti{Noi} and \iti{voi} \textbf{never} take the -isc- insertion.

Common: \iti{finire, preferire, pulire}

\chapter{Common Irregular Verbs} \label{M2-L3}
\begin{table}[ht]
    \centering
    \begin{tabular}{r | c c c c c c }
        con
        tent...
        Verb & io & tu & lui/lei & noi & voi & loro \\
        \hline
        \iti{andare} & \iti{vado} & \iti{vai} & \iti{va} & \iti{andiamo} & \iti{andate} & \iti{vanno}         \\
        \iti{fare}   & \iti{faccio} & \iti{fai} & \iti{fa} & \iti{facciamo} & \iti{fate} & \iti{fanno}        \\
        \iti{venire} & \iti{vengo} & \iti{vieni} & \iti{viene} & \iti{veniamo} & \iti{venite} & \iti{vengono} \\
        \iti{uscire} & \iti{esco} & \iti{esci} & \iti{esce} & \iti{usciamo} & \iti{uscite} & \iti{escono}     \\
        \iti{dare}   & \iti{do} & \iti{dai} & \iti{dà} & \iti{diamo} & \iti{date} & \iti{danno}               \\
        \iti{stare}  & \iti{sto} & \iti{stai} & \iti{sta} & \iti{stiamo} & \iti{state} & \iti{stanno}         \\
    \end{tabular}
\end{table}

\textbf{Key patterns:}
\begin{itemize}[nosep]
    \item \iti{-go / -gono} pattern: \iti{vengo/vengono, esco/escono}
    \item \iti{Noi} and \iti{voi} are usually the "calm" regular-looking forms
\end{itemize}

\chapter{Present Progressive — \iti{Stare} + Gerund} \label{M2-L4}

Used for actions happening \textbf{right now}: \iti{sto mangiando} — I am eating.

\section{Forming the Gerund}
\begin{table}[ht]
    \centering
    \begin{tabular}{r l c}
        Verb family & Gerund ending & Example \\
        \hline
        -ARE & \itb{-ando} & \iti{parlare} → \iti{parlando} \\
        -ERE & \itb{-endo} & \iti{leggere} → \iti{leggendo} \\
        -IRE & \itb{-endo} & \iti{dormire} → \iti{dormendo} \\
    \end{tabular}
\end{table}


\section{Irregular Gerunds}
\begin{itemize}[nosep]
    \item \iti{fare} → \itb{facendo}
    \item \iti{bere} → \itb{bevendo}
    \item \iti{dire} → \itb{dicendo}
\end{itemize}

\section{Structure}
\textbf{conjugated \iti{stare}} + gerund
\begin{itemize}[nosep]
    \item \iti{Sto mangiando.} — I am eating.
    \item \iti{Sta dormendo.} — He/she is sleeping.
    \item \iti{Stanno scrivendo.} — They are writing.
\end{itemize}

Italian uses the present progressive \textbf{less} than English — simple present is
often more natural for habits and near-future plans. \iti{Andare} almost never appears in the progressive.

%----------------------------------------------------------------------------------------------------
% END OF CHAPTER
%----------------------------------------------------------------------------------------------------